\documentclass[12pt]{article}
\usepackage[T1]{fontenc}
\usepackage[utf8]{inputenc}
\usepackage{lmodern}
\usepackage{textcomp}
\usepackage{enumitem}
\usepackage{geometry}
\usepackage{graphicx}
\usepackage{amsmath, amssymb}
\geometry{margin=1in}
\begin{document}
\begin{center}
Chemistry - Graduate Level Assessment 1 \
Total Marks: 40 \
Answer all questions.
\end{center}

\section*{Multiple Choice (10 marks)}
\begin{enumerate}[label=\arabic*.]
\item Silver bromide, AgBr, used in photography, may be prepared from AgNO\_3 and NaBr. Calculate
the weight of each required for producing 93.3 lb of AgBr. (1 lb = 454 g.)
\begin{enumerate}[label=(\alph*)]
    \item 50 lbs of NaBr and 85 lbs of AgNO3
    \item 45 lbs of NaBr and 75 lbs of AgNO3
    \item 52.7 lbs of NaBr and 84.7 lbs of AgNO3
    \item 70 lbs of NaBr and 50 lbs of AgNO3
    \item 40 lbs of NaBr and 90 lbs of AgNO3
\end{enumerate}
\item A buffer is made from equal concentrations of a weak acid and its conjugate base. Doubling
the volume of the buffer solution by adding water has what effect on its pH?
\begin{enumerate}[label=(\alph*)]
    \item It causes the pH to fluctuate unpredictably.
    \item It causes the pH to equalize with the water's pH.
    \item It causes the pH to become extremely basic.
    \item It neutralizes the pH completely.
    \item It significantly decreases the pH.
\end{enumerate}
\item The molecule with a tetrahedral shape is
\begin{enumerate}[label=(\alph*)]
    \item SO2
    \item BeCl2
    \item SF6
    \item CBr4
    \item PCl4F
\end{enumerate}
\item The Pauli exclusion principle states that
\begin{enumerate}[label=(\alph*)]
    \item no two electrons can occupy separate orbitals
    \item no two protons can have the same four quantum numbers
    \item two electrons can share the same quantum number if they are in different orbitals
    \item no two electrons can have the same four quantum numbers
    \item no two electrons can pair up if there is an empty orbital available
\end{enumerate}
\item Calculate the percentage difference between $e^x$ and $1+x$ for $x=0.0050$
\begin{enumerate}[label=(\alph*)]
    \item 7.50 $10^{-3} \%$
    \item 2.50 $10^{-3} \%$
    \item 0.75 $10^{-3} \%$
    \item 2.00 $10^{-3} \%$
    \item 1.25 $10^{-3} \%$
\end{enumerate}
\end{enumerate}

\section*{True / False (10 marks)}
\textit{Answer True or False}
\begin{enumerate}[label=\arabic*.]
\setcounter{enumi}{5}
\item True or False: The molality of a solution containing 2.3 g of ethanol in 500 g of water is 0.1 molal.
\item True or False: The dielectric constant of a non-polar gas like SO\_2 remains constant regardless of temperature and pressure changes.
\item True or False: The missing mass of the fluorine-19 isotope, (\^{}19\_9)F, is approximately 0.1756 amu, which is close to the atomic weight.
\item True or False: The equation 6x\^{}2 - 7x - 20 = 0 has solutions x = 5 and x = -4.
\item True or False: The units for the rate of a chemical reaction are expressed as L mol-1 s-1.
\end{enumerate}

\section*{Long Form Response (20 marks)}
\textit{Answer each question with a clear and complete explanation.}
\begin{enumerate}[label=\arabic*.]
\setcounter{enumi}{10}
\item Calculate the total length of a single DNA molecule in an E. coli chromosome, given its
parameters. Discuss its significance in relation to the size of the E. coli cell.
\item Discuss the significance of the intensity ratios in the EPR spectrum of the t-Bu radical
(CH3)3C•, and analyze the factors that influence these ratios in radical species.
\end{enumerate}

\vfill
\noindent\textit{End of Paper}
\end{document}